\section{Khảo sát thực trạng quy trình sử dụng thuốc nội trú (Mô hình As-Is)}
\label{sec:clinical_workflow_survey}

Qua khảo sát thực tế tại các khoa lâm sàng (Nội tổng hợp, Hồi sức tích cực, Tim mạch) tại các bệnh viện tuyến tỉnh và tuyến trung ương, quy trình kê đơn và cấp phát thuốc hiện tại (Mô hình As-Is) diễn ra qua các bước chính sau:

\begin{enumerate}[leftmargin=2em]
    \item \textbf{Thăm khám và chỉ định thuốc:} Bác sĩ điều trị thăm khám người bệnh tại giường, đối chiếu các kết quả xét nghiệm cận lâm sàng (creatinine, men gan, điện giải đồ) và nhập y lệnh thuốc vào phần mềm Quản lý Bệnh viện (HIS) hoặc viết tay vào sổ y lệnh.
    \item \textbf{Tra cứu thông tin dược học thủ công:} Khi gặp các ca bệnh phức tạp hoặc cần phối hợp từ 5 thuốc trở lên, bác sĩ phải tra cứu tương tác thuốc và liều lượng bằng cách lật sách Dược thư Quốc gia bản in, tra cứu tờ hướng dẫn sử dụng hoặc mở các ứng dụng điện thoại tra cứu độc lập.
    \item \textbf{Duyệt đơn và cấp phát dược lâm sàng:} Khoa Dược tiếp nhận đơn thuốc tổng hợp từ các khoa phòng. Dược sĩ lâm sàng tiến hành rà soát đơn thuốc theo đợt. Nếu phát hiện tương tác nguy hiểm hoặc liều bất thường, dược sĩ phải liên hệ lại với bác sĩ điều trị qua điện thoại để đề xuất đổi thuốc.
    \item \textbf{Thực hiện y lệnh tại buồng bệnh:} Điều dưỡng tiếp nhận thuốc từ tủ trực hoặc kho dược, đối chiếu 5 đúng (Đúng người bệnh, đúng thuốc, đúng liều, đúng đường dùng, đúng thời gian) và tiêm/truyền cho người bệnh.
\end{enumerate}

\noindent\textbf{Các điểm nghẽn và rủi ro cốt tử trong quy trình hiện tại:}
\begin{itemize}[leftmargin=1.5em]
    \item \textbf{Độ trễ cảnh báo lớn:} Khâu kiểm tra tương tác tại Khoa Dược diễn ra sau khi đơn thuốc đã được tạo. Việc can thiệp muộn làm gián đoạn kế hoạch điều trị của người bệnh.
    \item \textbf{Quá tải nhận thức của bác sĩ:} Một bác sĩ trong ca trực phải xử lý 30--50 bệnh nhân nội trú. Việc ghi nhớ toàn bộ hơn 1.400 chuyên luận thuốc và hàng nghìn cặp tương tác là bất khả thi đối với con người.
    \item \textbf{Nguy cơ xuất toán BHYT cao:} Bác sĩ thường chỉ tập trung vào hiệu quả điều trị mà khó nắm hết các ràng buộc hành chính phân tuyến kỹ thuật của Thông tư 20/2022/TT-BYT, dẫn đến tỷ lệ kê thuốc vượt tuyến bị xuất toán cao.
\end{itemize}
